\section*{附 录1 主要符号、缩写说明、方法对比表格}
\addcontentsline{toc}{section}{附 录1 主要符号、缩写说明、方法对比表格}

\begin{longtable}[]{@{}
  >{\raggedright\arraybackslash}p{(\columnwidth - 2\tabcolsep) * \real{0.1551}}
  >{\raggedright\arraybackslash}p{(\columnwidth - 2\tabcolsep) * \real{0.8449}}@{}}
\toprule\noalign{}
\begin{minipage}[b]{\linewidth}\raggedright
\textbf{缩写/符号}
\end{minipage} & \begin{minipage}[b]{\linewidth}\raggedright
\textbf{含义}
\end{minipage} \\
\midrule\noalign{}
\endhead
\bottomrule\noalign{}
\endlastfoot
ROS & Robot Operating System，机器人操作系统 \\
Gazebo & 机器人三维物理仿真平台 \\
URDF & Unified Robot Description Format，机器人统一描述格式 \\
A-SMGCS & 高级场面活动引导与控制系统 \\
LiDAR & 激光雷达 \\
\(v\) & 车辆纵向速度 \\
\(\omega\) & 车辆偏航角速度 \\
\(\delta_{f}\) & 前桥转角 \\
\(\delta_{r}\) & 后桥转角 \\
\(e_{x}\) & 纵向误差 \\
\(e_{y}\) & 横向误差 \\
\(e_{\theta}\) & 姿态误差 \\
\end{longtable}

方法对比表格

\begin{longtable}[]{@{}
  >{\raggedright\arraybackslash}p{(\columnwidth - 8\tabcolsep) * \real{0.1297}}
  >{\raggedright\arraybackslash}p{(\columnwidth - 8\tabcolsep) * \real{0.2442}}
  >{\raggedright\arraybackslash}p{(\columnwidth - 8\tabcolsep) * \real{0.2060}}
  >{\raggedright\arraybackslash}p{(\columnwidth - 8\tabcolsep) * \real{0.1866}}
  >{\raggedright\arraybackslash}p{(\columnwidth - 8\tabcolsep) * \real{0.2336}}@{}}
\toprule\noalign{}
\begin{minipage}[b]{\linewidth}\raggedright
\textbf{对比对象}
\end{minipage} & \begin{minipage}[b]{\linewidth}\raggedright
\textbf{主要特点}
\end{minipage} & \begin{minipage}[b]{\linewidth}\raggedright
\textbf{优点}
\end{minipage} & \begin{minipage}[b]{\linewidth}\raggedright
\textbf{局限性}
\end{minipage} & \begin{minipage}[b]{\linewidth}\raggedright
\textbf{与本文方法的关系}
\end{minipage} \\
\midrule\noalign{}
\endhead
\bottomrule\noalign{}
\endlastfoot
人工牵引 & 驾驶员和地面人员协同完成牵引作业 &
经验判断灵活，适应真实复杂环境能力强 &
依赖人员经验，标准化和重复精度受人为因素影响 &
本文方法不能直接替代人工，但可为自动化流程提供仿真验证基础 \\
固定路径/开环脚本 & 通过预设速度、转向角和时间完成动作 &
实现简单，适合早期动画和模型验证 &
缺少误差反馈，环境或初始状态变化后适应性差 &
本文在此基础上增加相对位姿误差控制思想 \\
PID控制 & 对速度、角度等单变量进行反馈调节 & 结构简单，工程应用广 &
难以单独处理多变量耦合的精确对接问题 &
可作为底层控制基础，但不足以完整描述对接策略 \\
Pure Pursuit & 基于前视点进行路径跟踪 & 路径跟踪平滑，实现难度适中 &
末端精确姿态对准能力有限 & 适合普通循迹，本文更强调近距离位姿误差修正 \\
Stanley算法 & 基于横向误差和航向误差进行控制 &
适合车辆路径跟踪，横向误差修正能力较好 &
更偏向中高速路径跟踪，低速精细对接针对性不足 &
可作为参考，但本文场景更偏低速对接 \\
MPC控制 & 在约束条件下进行预测优化控制 & 理论效果好，可处理多约束问题 &
建模和调参复杂，实现难度较高 & 后续可作为优化方向，本文阶段暂不采用 \\
SLAM/深度学习方法 & 依赖建图、视觉识别或数据驱动决策 &
自主性和智能化程度高 & 数据、算力和系统复杂度要求高 &
本文保留扩展接口，但现阶段以规则化仿真验证为主 \\
本文方法 & ROS-Gazebo仿真、飞机辅助感知、相对位姿误差控制、4WS执行 &
流程完整、结构清晰、实现难度适中，适合毕业设计验证 &
仿真简化较多，真实环境鲁棒性仍需进一步验证 &
作为自动牵引对接系统的前期仿真与算法验证方案 \\
\end{longtable}